\section{Previous works}

\paragraph{Software dependency networks.}
Dependency graphs arise naturally in software ecosystems: package managers such as npm, PyPI, and Maven define directed acyclic graphs in which each package declares its requirements~\cite{decan2019,labelle2004}. Studies of these ecosystems have revealed heavy-tailed degree distributions, small-world connectivity, and fragility under targeted removal of hub packages. Our analysis applies a similar methodology to a formal mathematics library---a setting in which every edge carries a stronger guarantee than ``package~$A$ requires package~$B$ at build time.'' In a type-checked library, every dependency has been verified by the kernel to be logically necessary; the graph records not merely a build requirement but a semantic relationship.

\paragraph{\module{Mathlib} growth and maintenance.}
Van Doorn, Ebner, and Lewis~\cite{vandoorn2020} described the engineering practices that sustain \module{Mathlib} as a collaboratively maintained library, addressing linting, documentation, and refactoring workflows. Commelin and Topaz~\cite{commelin2023} proposed abstraction boundaries and specification-driven development as organizing principles for formal mathematics. More recently, Commelin et al.~\cite{baanen2025} documented the challenges of scaling \module{Mathlib}, including build times, namespace reorganization, and technical debt. On the tooling side, Massot's \textsc{leanblueprint} system~\cite{massot2020leanblueprint} provides infrastructure for managing large formalization projects: mathematicians manually annotate dependency relations between statements in \LaTeX{} using \verb|\uses{}| tags, producing a dependency graph that tracks formalization progress. This approach has been adopted by several major projects, including the Liquid Tensor Experiment and the ongoing formalization of Fermat's Last Theorem. Zhu et al.~\cite{zhu2025leanarchitect} complement this with \textsc{LeanArchitect}, which automatically infers dependency information from Lean source code, eliminating the manual duplication between informal and formal representations. Notably, their work includes AI integration: in a case study, the automatically extracted dependency graph guides an AI prover (Achim et al.'s Aristotle system~\cite{achim2025aristotle}) to formalize theorems in topological order. Our work differs from both in scope and purpose: \textsc{leanblueprint} and \textsc{LeanArchitect} extract dependency graphs to manage individual formalization projects; we analyze the global dependency structure of the entire \module{Mathlib} library as a network, applying tools from network science---degree distributions, community detection, robustness analysis---to characterize structural patterns that no single project would reveal. These works address \module{Mathlib}'s maintenance from a software engineering perspective. Our analysis complements theirs by providing a network-theoretic diagnostic of the same challenges they address qualitatively: we show, for example, that the $17.5\%$ import redundancy rate and $0.107$ module cohesion quantify the cognitive overhead that maintenance practices must manage.

\paragraph{Machine learning for formal mathematics.}
The graph structure of \module{Mathlib} has been exploited in machine learning, particularly for premise selection. Yang et al.~\cite{leandojo} introduced the LeanDojo benchmark, extracting proof traces from \module{Mathlib} that implicitly define a dependency graph. Bemberek et al.~\cite{bemberek2025} constructed a heterogeneous dependency graph from this benchmark and applied relational graph convolutional networks for premise selection, achieving a $25\%$ improvement over text-only baselines. Lu et al.~\cite{leanfinder} built a directed acyclic graph of Lean~4 statements with dependency edges to support semantic search. On the knowledge-representation side, the EpiK Protocol~\cite{epik2024} constructed a knowledge graph from Lean~4 and \module{Mathlib}, applying embedding models for link prediction. Uskuplu and Moss~\cite{knowtex2025} argued that dependency graphs should become a standard feature of mathematical writing. These systems consume the dependency graph as training data or as a retrieval index. We study the graph itself---not as input to a downstream task, but as an object of independent structural interest.

\paragraph{Complex network methodology.}
The analytical tools we employ---degree distributions, centrality measures, community detection---originate in network science. Heavy-tailed network theory~\cite{barabasi1999}, small-world analysis~\cite{watts1998}, and modularity-based community detection~\cite{newman2004,blondel2008} have been applied to biological, social, and technological networks, but not, to our knowledge, systematically to a large-scale formal mathematics library. In a related setting, Barbosa et al.~\cite{barbosa2013} studied the network structure of three online mathematical libraries---Wikipedia (restricted to mathematics), MathWorld, and DLMF---finding large strongly connected components, the absence of clear power laws, and the effectiveness of stress centrality for navigating mathematical content. Their objects of study, however, are hyperlinked encyclopedias with informal content; ours is a machine-verified library in which every dependency is logically certified.

\paragraph{Gap.}
Prior work uses the dependency graph instrumentally---as training data for machine learning, as a compilation dependency for build systems, or as a navigational aid for human developers. No study has analyzed the graph itself as a structural object encoding both inherited mathematical structure and human organizational process. This paper takes a first step toward filling that gap.
